\subsection{Polynomial rings over a Krull domain}

PROPOSITION 13. Let A be a Krull domain and \(\mathbf{X}_1, \mathbf{X}_2, \ldots, \mathbf{X}_n\) indeterminates. The ring \(\mathrm{A}[\mathbf{X}_1, \ldots, \mathbf{X}_n]\) is a Krull domain.

Arguing by induction on n, it is sufficient to show that, if X is an indeterminate, A{[}X{]} is a Krull domain. Let K be the field of fractions of A. The field of fractions of A{[}X{]} is K(X). Let I be the set of monic polynomials in K{[}X{]} which are irreducible over K; for all \(f \in I\) , let \(v_{f}\) be the valuation on K(X) defined by f (Chapter VI, §3, no. 3, Example 4). On the other hand, for every essential valuation w of A, let \(\bar{w}\) be the extension of w to K(X) defined by

\[
\bar {w} \left(\sum_ {j} a _ {j} X ^ {j}\right) = \inf _ {j} (w (a _ {j}))
\]

for \(\sum_{j} a_j X^j \in K[X]\) (Chapter VI, § 10, no. 1, Lemma 1). Clearly the \(v_f\) and the \(\bar{w}\) are discrete and normed and, for all \(u \in K[X]\) , \(v_f(u) = 0\) (resp. \(\bar{w}(u) = 0\) ) except for a finite number of valuations \(v_f\) (resp. \(\bar{w}\) ).

To show the proposition, it therefore suffices to show that A{[}X{]} is the intersection of the rings of the valuations \(v_{f}\) and \(\bar{w}\) . Now the intersection of the rings of the valuations \(v_{f}\) is K{[}X{]}. On the other hand, for \(\sum_{j} a_{j} X^{j} \in K[X]\) , the relation \(\bar{w}\left(\sum_{j} a_{j} X^{j}\right) \geqslant 0\) is equivalent to `` \(w(a_{j}) \geqslant 0\) for all j''; hence the relation `` \(\bar{w}\left(\sum_{j} a_{j} X^{j}\right) \geqslant 0\) for every valuation \(\bar{w}\) '' is equivalent to `` \(w(a_{j}) \geqslant 0\) for all j and every essential valuation w of A.'' This proves our assertion.

Remark. The valuations \(v_f\) and \(\bar{w}\) introduced in the proof of Proposition 13 are the essential valuations of A{[}X{]}. It will be sufficient for us to show that, if V is the set of valuations \(v_f\) ( \(f\) irreducible) and \(\bar{w}\) ( \(w\) essential), then, for all \(v' \in V\) , there exists an element \(g \in K(X)\) which is not in A{[}X{]} and such that \(v''(g) \geqslant 0\) for all the valuations \(v'' \in V\) distinct from \(v'\) ; this will prove that \(V - \{v'\}\) does not satisfy (AK \(_{\text{II}}\) ) and the conclusion will follow therefore from no. 4, Corollary 2 to Proposition 6. Suppose first that \(v'\) is of the form \(\bar{w}\) : then we may take \(g\) to be an element \(b \in K\) such that \(w(b) < 0\) , \(w'(b) \geqslant 0\) for the essential valuations \(w'\) of A distinct from \(w\) , for then \(v_f(b) = 0\) for every irreducible monic polynomial \(f\) in K{[}X{]}; the existence of an element \(b\) satisfying the above conditions follows from no. 5, Proposition 9. Suppose secondly that \(v'\) is of the form \(v_f\) for an irreducible monic polynomial \(f \in K[X]\) of degree \(m\) ; then we may take \(g = a/f\) where \(a \in A\) . For \(v_h(g) \geqslant 0\) for every irreducible monic polynomial \(h \neq f\) in K{[}X{]}; it remains to choose \(a \in A\) such that, for every essential valuation \(w\) of A, \(w(a)\) is at least equal to the greatest lower bound of the elements \(w(c_i)\) , where the \(c_i\) are the coefficients of \(f (1 \leqslant i \leqslant m)\) ; now the existence of such an \(a \in A\) follows from (AK \(_{\text{III}}\) ) and no. 5, Proposition 9.

We may also say (no. 6, Theorem 4) that the prime ideals of A{[}X{]} of height 1 are:

\begin{enumerate}
\def\labelenumi{(\arabic{enumi})}
\item
  the prime ideals of the form \(\mathfrak{p}\mathbf{A}[\mathbf{X}]\) , where \(\mathfrak{p}\) is a prime ideal of \(\mathbf{A}\) of height 1;
\item
  the prime ideals of the form \(\mathfrak{m} \cap \mathrm{A}[\mathrm{X}]\) , where \(\mathfrak{m}\) is a (necessarily principal) prime ideal of \(\mathrm{K}[\mathrm{X}]\) .
\end{enumerate}

The latter are characterized by the fact that their intersection with A is reduced to 0.

